\scanpage{180}
\mspage{129}
Se donner une flèche dans $\mathit{Cat}$
\[
\mathscr X\longrightarrow Y
\]
\textquotedblleft revient au $=$\textquotedblright\ que de se donner un système $(f_0,f_1,f_2,\lambda_1,\lambda_2)$ comme dessus, \inserted{et ceci} \textquotedblleft fonctoriellement en $Y$\textquotedblright. C'est là \deleted{\unsure{le|les}} \unsure{cas|un} particulier de la détermination générale des foncteurs d'\deleted{une intégrale} la catégorie fibrée \textquotedblleft \unsure{intégrée|intégrale}\textquotedblright\ d'un \inserted{\unsure{contrav.}} foncteur
\[
I^0\longrightarrow\mathit{Cat},\qquad
\begin{aligned}i&\longmapsto X_i\\u&\longmapsto u^*\end{aligned}
\]
\inserted{vers une cat. $Y$,} en termes d'un système de foncteurs
\[
X_i\xrightarrow{f_i}Y ,
\]
et d'homomorphismes fonctoriels
\[
f_i\circ u^*\xrightarrow{\lambda_u} f_j
\]
pour $u:i\to j$ d'où $\unsure{u'|u^*}:X_j\to X_i$, avec des conditions de transitivité pour des flèches composées, dont on n'a pas à tenir compte dans le cas actuel, vu qu'il n'y a pas de flèches composées à regarder dans \deleted{$\unsure{\Phi}$} $I=\Phi$ (sauf des composés \deleted{\unsure{avec}} avec flèches identiques, \deleted{\unsure{quitte à}} qu'on peut se dispenser de \uline{regarder} ici \ldots).
% archive sheet number: 180
